\Subhead{6.4.1}{Стохастические модели сообщения и шума. Уравнения фильтра и Риккати. Сравнение с методом Колмогорова--Винера}

В непрерывном времени белый измерительный шум корректно записывать через приращения наблюдения, а не как обычную дифференцируемую функцию:
\begin{equation}
 \dd\bm x=A(t)\bm x\dd t+G(t)\dd\bm W,\qquad
 \dd\bm y=H(t)\bm x\dd t+\dd\bm V.
\end{equation}
Пусть независимые винеровские приращения имеют ковариации $Q\dd t$ и $R\dd t$, $R>0$, а начальное состояние гауссовское и независимо от шумов. Тогда апостериорное распределение остаётся гауссовским, достаточно среднего $\hat{\bm x}$ и ковариации ошибки $P$:
\begin{align}\label{eq:kalmanbucy}
 \dd\hat{\bm x}&=A\hat{\bm x}\dd t+K(\dd\bm y-H\hat{\bm x}\dd t),\quad K=PH^{\mathsf T}R^{-1},\\
 \dot P&=AP+PA^{\mathsf T}+GQG^{\mathsf T}-PH^{\mathsf T}R^{-1}HP.
\end{align}
Первый член прогноза следует модели, инновация сравнивает наблюдение с прогнозом. Ковариация растёт из-за динамического шума и уменьшается от информативного измерения.

Краткий вывод усиления: у ошибки $\bm e=\bm x-\hat{\bm x}$ дрейф $A-KH$, шумы $G\dd W-K\dd V$. Поэтому
\begin{equation}
 \dot P=(A-KH)P+P(A-KH)^{\mathsf T}+GQG^{\mathsf T}+KRK^{\mathsf T}.
\end{equation}
Дополнение квадрата по $K$ даёт минимум при $K=PH^{\mathsf T}R^{-1}$ и отрицательный квадратичный член Риккати. Начальные $\hat x_0,P_0$ — априорные среднее и неопределённость, а не произвольные нули.

Если под «апостериорной информацией» понимать матрицу точности $J=P^{-1}$, то $\dot J=-J\dot PJ$:
\begin{equation}
 \dot J=-JA-A^{\mathsf T}J-JGQG^{\mathsf T}J+H^{\mathsf T}R^{-1}H.
\end{equation}
Это также уравнение типа Риккати. Матрица точности не тождественна энтропии или количеству информации Шеннона.

Пример: $\dd x=-ax\dd t+\sqrt q\dd W$, $\dd y=hx\dd t+\sqrt r\dd V$. Тогда
\begin{equation}
 \dot P=q-2aP-h^2P^2/r,\quad
 P_\infty=\frac r{h^2}\left(\sqrt{a^2+h^2q/r}-a\right)
\end{equation}
при $h\ne0$, $q,r>0$, $a>0$. При ухудшении измерения $r\to\infty$ ошибка стремится к априорной дисперсии $q/(2a)$. При хороших измерениях уменьшается; не равна нулю при конечном шуме и непрерывно возбуждаемой динамике.

Винер работает с корреляциями/спектрами стационарной модели; Калман–Бьюси — с моделью состояния, допускает нестационарные коэффициенты и рекурсивное обновление. Для одной линейно-гауссовской стационарной модели после установления оба дают тот же причинный оптимум. Наблюдаемость, обнаруживаемость и условия возбуждения определяют сходимость Риккати; её нельзя обещать для любой матрицы $A,H$.
